% !TEX root = ../main.tex
\chapter{Туршилтын арга}\label{ch:method}

\section{Орчин}
Тестийг багийн гишүүний компьютер (macOS 27, Apple Silicon, 10 цөм, Eclipse Temurin 17.0.21) болон хөгжүүлэлтийн
Linux контейнер (Ubuntu 24.04, Intel Xeon, 2 цөм, OpenJDK 21) дээр ажиллуулав. Хугацааны хэмжилтийг контейнер дээр хийсэн.

\section{Автомат тест}
Тестийг JUnit-гүй, \code{tests/TestRunner.java} файлд бичсэн бөгөөд \code{./test.sh} ажиллуулна. Тест бүрд
урьдчилан бичсэн хүлээгдсэн үр дүн ба програмын бодит үр дүнг харьцуулж, бүх эрэмбэлэлтийг \code{Arrays.sort}-той
тулгана. Үр дүнгийн хүснэгт \code{tests/results/test-results.md}-д бичигдэнэ.

\begin{table}[H]
\centering\small
\caption{Автомат тестийн ангилал}\label{tab:tests}
\begin{tabularx}{\linewidth}{@{}lrY@{}}
\toprule
Ангилал & Тоо & Шалгах зүйл \\
\midrule
Номын жишээ & 26 & 25 алгоритмын хослол, mid = 5, хоёр хэсэг \\
Захын тохиолдол & 17 & $n = 1, 2, 3$, бүгд ижил, сөрөг, INT\_MIN/MAX, эрэмбэтэй, урвуу, $\lceil n/2\rceil$ / $\lfloor n/2\rfloor$ \\
Буруу оролт & 18 & Хоосон, тоо биш, бутархай, int-ээс гадуур, хязгаарууд, буруу файл \\
Давхардал ба хэлбэр & 6 & 4 хэлбэр × 25 хослол, 2000 санамсаргүй ажиллалт \\
Том өгөгдөл & 4 & $n = 10^6$ (merge, quick), $n = 20\,000$ (O($n^2$)) \\
Thread & 5 & 3 өөр thread, параметр, \code{TERMINATED}, дараалал, муж зөрчөөгүй \\
Event trace & 7 & Trace-ээс сэргээсэн \code{result[]} = бодит, JSON формат \\
Lockstep & 5 & \code{WAITING}, 1 permit = 1 үйлдэл, ээлжлэх, цуцлах \\
Командын мөр & 12 & Тусдаа JVM дээр гаралт ба exit code \\
\bottomrule
\end{tabularx}
\end{table}

Тест нь зөвхөн үр дүнг биш, номын дотоод нөхцлийг ч шалгадаг. Trace-ийн бүх event-ээр Sorter-L нь
\code{[0, mid)}-ээс, Sorter-R нь \code{[mid, n)}-ээс гадуур хандаагүйг баталсан. Мөн merger хоёр sorter дууссаны дараа
эхэлснийг 100 удаа шалгасан. Lockstep-ийн тестийг 30 удаа дараалан болон нэг цөм дээр (\code{taskset -c 0}) ажиллуулсан.
Web UI-г Playwright (Chromium)-оор автоматаар шалгаж, JavaScript алдаа 0 байгааг баталсан.

\section{Хугацааны хэмжилтийн арга}
\code{tests/Benchmark.java} (\code{./bench.sh}) нь гурван хувилбарыг харьцуулна: \emph{1 thread} (нэг thread бүх массивыг
эрэмбэлнэ), \emph{1 thread, хувааж} (номын бүтэцтэй ижил ажлыг нэг thread дараалан хийнэ) ба \emph{3 thread} (номын
бүтэц). Нийт хурдсалт $= T_{1}/T_{3}$, зэрэгцээний хурдсалт $= T_{1,\text{хувааж}}/T_{3}$. Хоёр дахь нь хуваалтын
нөлөөг зэрэгцээний нөлөөнөөс салгана. Оролт санамсаргүй \code{int}, trace унтраалттай. JIT-ийн нөлөөг багасгахын тулд
3 тусдаа JVM ажиллуулж, тус бүрд 3 удаа дулаацуулсны дараа 9 удаа хэмжиж, медианы медианыг авсан.
